\documentclass[12pt,a4paper]{article}
\usepackage[utf8]{inputenc}
\usepackage[T1]{fontenc}
\usepackage{amsmath,amssymb,amsthm}
\usepackage{graphicx}
\usepackage{listings}
\usepackage{xcolor}
\usepackage{hyperref}
\usepackage{geometry}
\usepackage{tcolorbox}
\usepackage{needspace}
\usepackage{tikz}
\usepackage{array}
\usepackage{booktabs}
\usepackage{pifont}
\usetikzlibrary{arrows,positioning,matrix}
\geometry{margin=2.5cm}

% Command for circled numbers
\newcommand*\circled[1]{\tikz[baseline=(char.base)]{\node[shape=circle,draw,inner sep=2pt] (char) {#1};}}

% Définition des couleurs personnalisées
\definecolor{titrebleu}{RGB}{0,102,204}
\definecolor{defvert}{RGB}{0,153,76}
\definecolor{exempleorange}{RGB}{255,102,0}
\definecolor{algoviolet}{RGB}{153,0,153}
\definecolor{importantrouge}{RGB}{204,0,0}
\definecolor{fondgris}{RGB}{245,245,245}
\definecolor{fondjaune}{RGB}{255,250,205}
\definecolor{fondvert}{RGB}{230,255,230}
\definecolor{fondbleu}{RGB}{230,240,255}
\definecolor{fondrose}{RGB}{255,230,240}

% Boîtes colorées
\newtcolorbox{definition}{
    colback=fondvert,
    colframe=defvert,
    fonttitle=\bfseries,
    title=Definition,
    arc=3mm
}

\newtcolorbox{exemple}{
    colback=fondjaune,
    colframe=exempleorange,
    fonttitle=\bfseries,
    title=Exemple,
    arc=3mm
}

\newtcolorbox{important}{
    colback=white,
    colframe=importantrouge,
    fonttitle=\bfseries,
    title=Important,
    arc=3mm
}

\newtcolorbox{astuce}{
    colback=fondbleu,
    colframe=titrebleu,
    fonttitle=\bfseries,
    title=A savoir,
    arc=3mm
}

\newtcolorbox{methode}{
    colback=fondbleu,
    colframe=algoviolet,
    fonttitle=\bfseries,
    title=Methode,
    arc=3mm
}

\newtcolorbox{notecours}{
    colback=fondrose,
    colframe=importantrouge,
    fonttitle=\bfseries,
    title=Note de Cours,
    arc=3mm
}

% Style des sections
\usepackage{titlesec}
\titleformat{\section}
{\color{titrebleu}\normalfont\Large\bfseries}
{\color{titrebleu}\thesection}{1em}{}

\titleformat{\subsection}
{\color{defvert}\normalfont\large\bfseries}
{\color{defvert}\thesubsection}{1em}{}

\title{\textbf{\Huge\color{titrebleu}La Methode du Simplexe}\\[0.5em]
\large\color{defvert}Comment Fonctionne la Methode du Simplexe ?\\
\small Modeles Lineaires de la Recherche Operationnelle}
\author{\large Prof. Herve KERIVIN, PhD\\
\small Universite Clermont Auvergne\\
Clermont Auvergne INP - ISIMA\\
LIMOS UMR 6158 CNRS}
\date{\large Version 2024-2025 - Cours 2}

\begin{document}

\maketitle
\tableofcontents
\newpage

\section{Introduction}

\needspace{10\baselineskip}
\subsection{Historique}

\begin{notecours}
\textbf{Naissance de la programmation lineaire :}

La programmation lineaire est relativement \textcolor{titrebleu}{\textbf{recente}}.

Elle a commence en \textcolor{importantrouge}{\textbf{1947}} lorsque \textcolor{defvert}{\textbf{George B. Dantzig}} a concu la \textbf{``Methode du Simplexe''} pour resoudre des problemes de planification de la Force Aerienne americaine (U.S. Air Force) formules en programmation lineaire.
\end{notecours}

\begin{astuce}
\textbf{Importance historique :}

\begin{itemize}
    \item La methode du simplexe est l'un des algorithmes les plus utilises et les plus importants du 20e siecle
    \item Elle reste la methode de reference pour resoudre des problemes de PL en pratique
    \item Bien qu'elle ne soit pas polynomiale dans le pire cas, elle est extremement efficace en pratique
\end{itemize}
\end{astuce}

\newpage
\section{Motivation : Un Premier Exemple}

\needspace{10\baselineskip}
\subsection{Enonce du Probleme}

\begin{exemple}
\textbf{Considerons le programme lineaire suivant :}

\begin{align*}
\text{maximiser} \quad & z = -x_1 - 2x_2\\
\text{sujet a} \quad & x_1 + x_2 \leq 3\\
& 3x_1 - x_2 \leq 5\\
& x_1 \geq 0\\
& x_2 \geq 0
\end{align*}
\end{exemple}

\begin{astuce}
\textbf{Solution evidente :}

\begin{itemize}
    \item \textcolor{defvert}{\textbf{Solution realisable evidente :}} $x_1 = 0$, $x_2 = 0$
    \item \textcolor{titrebleu}{\textbf{Cette solution est optimale !}}
    \item \textbf{Pourquoi ?} Augmenter soit $x_1$ soit $x_2$ diminuerait la valeur de $z$
\end{itemize}

\textbf{Valeur de l'objectif :} $z = 0$
\end{astuce}

\begin{notecours}
\textbf{Observation importante :}

Dans cet exemple simple, la solution a l'origine $(0,0)$ est optimale car les coefficients de la fonction objectif sont tous \textcolor{importantrouge}{\textbf{negatifs}}.

Mais comment proceder dans le cas general ?
\end{notecours}

\newpage
\section{Variables d'Ecart (Slack Variables)}

\needspace{10\baselineskip}
\subsection{Concept de Variable d'Ecart}

\begin{definition}
\textbf{Variable d'ecart (Slack variable) :}

Considerons l'inegalite $x_1 + x_2 \leq 3$.

Pour toute solution realisable $x_1, x_2$, le membre gauche de cette contrainte est au plus egal a la valeur du membre droit ; souvent, il peut y avoir un \textcolor{defvert}{\textbf{ecart}} entre les deux valeurs.

Notons $x_3$ l'ecart, c'est-a-dire :
$$x_3 = 3 - x_1 - x_2$$

Avec cette notation, l'inegalite $x_1 + x_2 \leq 3$ peut maintenant etre ecrite comme $x_3 \geq 0$.
\end{definition}

\begin{exemple}
\textbf{Exemples numeriques :}

\begin{itemize}
    \item Si $x_1 = 1$ et $x_2 = 1$, alors $x_3 = 3 - 1 - 1 = 1$ (il reste un ecart de 1)
    \item Si $x_1 = 2$ et $x_2 = 1$, alors $x_3 = 3 - 2 - 1 = 0$ (la contrainte est \textit{saturee})
    \item Si $x_1 = 0$ et $x_2 = 0$, alors $x_3 = 3$ (ecart maximal)
\end{itemize}
\end{exemple}

\newpage
\subsection{Transformation en Systeme d'Equations}

\begin{methode}
\textbf{Transformation des contraintes :}

On transforme les contraintes d'un probleme de PL (ecrit sous forme standard) en un systeme d'equations lineaires en ajoutant des variables d'ecart.

\textbf{Forme originale :}
\begin{align*}
a_{11}x_1 + a_{12}x_2 + \cdots + a_{1n}x_n &\leq b_1\\
a_{21}x_1 + a_{22}x_2 + \cdots + a_{2n}x_n &\leq b_2\\
&\vdots\\
a_{m1}x_1 + a_{m2}x_2 + \cdots + a_{mn}x_n &\leq b_m\\
x_1, x_2, \ldots, x_n &\geq 0
\end{align*}

\textbf{Avec variables d'ecart :}
\begin{align*}
a_{11}x_1 + a_{12}x_2 + \cdots + a_{1n}x_n + x_{n+1} &= b_1\\
a_{21}x_1 + a_{22}x_2 + \cdots + a_{2n}x_n + x_{n+2} &= b_2\\
&\vdots\\
a_{m1}x_1 + a_{m2}x_2 + \cdots + a_{mn}x_n + x_{n+m} &= b_m\\
x_1, x_2, \ldots, x_n, x_{n+1}, \ldots, x_{n+m} &\geq 0
\end{align*}
\end{methode}

\begin{important}
\textbf{Remarque cruciale :}

Les variables d'ecart sont \textcolor{importantrouge}{\textbf{non negatives}} (c'est-a-dire $x_{n+1} \geq 0$, $x_{n+2} \geq 0$, \ldots, $x_{n+m} \geq 0$).
\end{important}

\newpage
\subsection{Premier Exemple avec Variables d'Ecart}

\begin{exemple}
\textbf{Programme lineaire :}

\begin{align*}
\text{maximiser} \quad & z = 5x_1 + 4x_2 + 3x_3\\
\text{sujet a} \quad & 2x_1 + 3x_2 + x_3 \leq 5\\
& 4x_1 + x_2 + 2x_3 \leq 11\\
& 3x_1 + 4x_2 + 2x_3 \leq 8\\
& x_1, x_2, x_3 \geq 0
\end{align*}

\textbf{Solutions realisables evidentes :}
\begin{itemize}
    \item $x_1 = 0, x_2 = 0, x_3 = 0$ est realisable
    \item $x_1 = \frac{5}{2}, x_2 = 0, x_3 = 0$ est realisable
\end{itemize}
\end{exemple}

\begin{exemple}
\textbf{Avec variables d'ecart :}

Definissons :
\begin{align*}
x_4 &= 5 - 2x_1 - 3x_2 - x_3\\
x_5 &= 11 - 4x_1 - x_2 - 2x_3\\
x_6 &= 8 - 3x_1 - 4x_2 - 2x_3\\
z &= 5x_1 + 4x_2 + 3x_3
\end{align*}

ou $x_4$, $x_5$ et $x_6$ sont les variables d'ecart pour les premiere, deuxieme et troisieme inegalites, respectivement.

Le probleme devient :
\begin{align*}
\text{maximiser} \quad & z = 5x_1 + 4x_2 + 3x_3\\
\text{sujet a} \quad & x_4 = 5 - 2x_1 - 3x_2 - x_3\\
& x_5 = 11 - 4x_1 - x_2 - 2x_3\\
& x_6 = 8 - 3x_1 - 4x_2 - 2x_3\\
& x_1, x_2, x_3, x_4, x_5, x_6 \geq 0
\end{align*}

\textbf{Solutions correspondantes :}
\begin{itemize}
    \item $x_1 = 0, x_2 = 0, x_3 = 0, x_4 = 5, x_5 = 11, x_6 = 8$
    \item $x_1 = \frac{5}{2}, x_2 = 0, x_3 = 0, x_4 = 0, x_5 = 1, x_6 = \frac{1}{2}$
\end{itemize}
\end{exemple}

\newpage
\section{Dictionnaires}

\needspace{10\baselineskip}
\subsection{Dictionnaire Initial}

\begin{definition}
\textbf{Dictionnaire initial :}

Etant donne un programme lineaire ecrit sous forme standard :

\begin{align*}
\text{maximiser} \quad & z = \sum_{j=1}^{n} c_jx_j\\
\text{sujet a} \quad & \sum_{j=1}^{n} a_{ij}x_j \leq b_i \quad \text{pour } i = 1, 2, \ldots, m\\
& x_j \geq 0 \quad \text{pour } j = 1, 2, \ldots, n
\end{align*}

On introduit les variables d'ecart $x_{n+1}, x_{n+2}, \ldots, x_{n+m} \geq 0$.

Le \textcolor{defvert}{\textbf{dictionnaire initial}} est :

\begin{align*}
x_{n+i} &= b_i - \sum_{j=1}^{n} a_{ij}x_j \quad \text{pour } i = 1, 2, \ldots, m\\
z &= \sum_{j=1}^{n} c_jx_j
\end{align*}
\end{definition}

\begin{exemple}
\textbf{Dictionnaire initial pour notre exemple :}

\begin{align*}
x_4 &= 5 - 2x_1 - 3x_2 - x_3\\
x_5 &= 11 - 4x_1 - x_2 - 2x_3\\
x_6 &= 8 - 3x_1 - 4x_2 - 2x_3\\
z &= 0 + 5x_1 + 4x_2 + 3x_3
\end{align*}
\end{exemple}

\newpage
\subsection{Variables de Base et Hors-Base}

\begin{definition}
\textbf{Variables de base et hors-base :}

Les equations de chaque dictionnaire doivent exprimer $m$ des variables $x_1, x_2, \ldots, x_{n+m}$ (appelees \textcolor{titrebleu}{\textbf{variables de base}}) et la fonction objectif $z$ en termes des $n$ variables restantes (appelees \textcolor{defvert}{\textbf{variables hors-base}}).

\textbf{Dans le dictionnaire initial :}
\begin{itemize}
    \item \textcolor{titrebleu}{\textbf{Variables de base :}} $x_4, x_5, x_6$ (variables d'ecart)
    \item \textcolor{defvert}{\textbf{Variables hors-base :}} $x_1, x_2, x_3$ (variables de decision)
\end{itemize}
\end{definition}

\begin{definition}
\textbf{Base :}

Les variables de base constituent une \textcolor{algoviolet}{\textbf{base}}.

\textbf{Dans le dictionnaire initial :} la base est $B = \{4, 5, 6\}$ (ou $\{x_4, x_5, x_6\}$).
\end{definition}

\begin{notecours}
\textbf{Notation importante :}

Pour un dictionnaire associe a une base $B$ :
\begin{itemize}
    \item $i \in B$ signifie que $x_i$ est une variable de base
    \item $j \notin B$ (ou $j \in N$) signifie que $x_j$ est une variable hors-base
    \item $N$ designe l'ensemble des indices des variables hors-base
\end{itemize}
\end{notecours}

\newpage
\subsection{Dictionnaire General}

\begin{definition}
\textbf{Dictionnaire associe a une base $B$ :}

Soit $B$ l'ensemble des indices des $m$ variables de base du PL. Le dictionnaire associe a $B$ est :

\begin{align*}
x_i &= \bar{b}_i - \sum_{j \notin B} \bar{a}_{ij}x_j \quad \text{pour tout } i \in B\\
z &= \bar{z} + \sum_{j \notin B} \bar{c}_j x_j
\end{align*}

ou :
\begin{itemize}
    \item $\bar{b}_i$ : valeur de la variable de base $x_i$ lorsque les variables hors-base sont nulles
    \item $\bar{a}_{ij}$ : coefficients des variables hors-base dans les equations
    \item $\bar{z}$ : valeur de l'objectif lorsque les variables hors-base sont nulles
    \item $\bar{c}_j$ : coefficients des variables hors-base dans l'objectif
\end{itemize}
\end{definition}

\begin{exemple}
\textbf{Exemple de dictionnaire different :}

Base $B = \{1, 3, 5\}$ :

\begin{align*}
x_1 &= 2 - 2x_2 - 2x_4 + x_6\\
x_3 &= 1 + x_2 + 3x_4 - 2x_6\\
x_5 &= 1 + 5x_2 + 2x_4\\
z &= 13 - 3x_2 - x_4 - x_6
\end{align*}

Variables hors-base : $x_2, x_4, x_6$
\end{exemple}

\newpage
\section{Dictionnaires Realisables}

\needspace{10\baselineskip}
\subsection{Definition}

\begin{definition}
\textbf{Dictionnaire realisable :}

Si, lorsqu'on :
\begin{enumerate}
    \item Fixe les $n$ variables du membre droit a zero, et
    \item Evalue les $m$ variables du membre gauche,
\end{enumerate}

on arrive a une solution realisable, alors le dictionnaire est appele un \textcolor{defvert}{\textbf{dictionnaire realisable}}.
\end{definition}

\begin{exemple}
\textbf{Verification de la realisabilite :}

\textbf{Dictionnaire initial :}
\begin{align*}
x_4 &= 5 - 2x_1 - 3x_2 - x_3\\
x_5 &= 11 - 4x_1 - x_2 - 2x_3\\
x_6 &= 8 - 3x_1 - 4x_2 - 2x_3\\
z &= 0 + 5x_1 + 4x_2 + 3x_3
\end{align*}

\textbf{Mise a zero des variables hors-base :} $x_1 = x_2 = x_3 = 0$

\textbf{Evaluation des variables de base :} $x_4 = 5$, $x_5 = 11$, $x_6 = 8$

\textbf{Toutes sont non negatives !} $\Rightarrow$ \textcolor{defvert}{\textbf{Dictionnaire realisable}}
\end{exemple}

\begin{exemple}
\textbf{Autre dictionnaire realisable :}

\begin{align*}
x_1 &= 2 - 2x_2 - 2x_4 + x_6\\
x_3 &= 1 + x_2 + 3x_4 - 2x_6\\
x_5 &= 1 + 5x_2 + 2x_4\\
z &= 13 - 3x_2 - x_4 - x_6
\end{align*}

\textbf{Variables hors-base a zero :} $x_2 = x_4 = x_6 = 0$

\textbf{Variables de base :} $x_1 = 2$, $x_3 = 1$, $x_5 = 1$

\textcolor{defvert}{\textbf{Toutes non negatives $\Rightarrow$ Realisable}}
\end{exemple}

\newpage
\subsection{Solutions de Base}

\begin{definition}
\textbf{Propositions importantes :}

\textbf{Proposition 1 :} Chaque dictionnaire realisable decrit une solution realisable.

\textbf{Proposition 2 :} Toute solution realisable n'est pas necessairement decrite par un dictionnaire realisable.
\end{definition}

\begin{exemple}
\textbf{Contre-exemple pour la Proposition 2 :}

Solution realisable : $x_1 = 1, x_2 = 0, x_3 = 1, x_4 = 2, x_5 = 5, x_6 = 3$

Cette solution a \textcolor{importantrouge}{\textbf{cinq variables non nulles}}.

Or, dans un dictionnaire realisable, on a au plus \textcolor{defvert}{\textbf{trois variables non nulles}} (les variables de base, car $m = 3$).

Donc cette solution \textcolor{importantrouge}{\textbf{ne peut pas}} etre decrite par un dictionnaire.
\end{exemple}

\begin{definition}
\textbf{Solutions de base :}

Les solutions realisables qui peuvent etre decrites par des dictionnaires sont appelees \textcolor{algoviolet}{\textbf{solutions de base}} (ou \textbf{solutions de base realisables}).
\end{definition}

\begin{astuce}
\textbf{Caracteristique d'une solution de base :}

Une solution de base a \textcolor{titrebleu}{\textbf{au plus $m$ variables non nulles}} (les variables de base).

Les $n$ autres variables (hors-base) sont nulles.
\end{astuce}

\newpage
\subsection{Origine Realisable}

\begin{definition}
\textbf{Origine et origine realisable :}

Dans le dictionnaire initial :
\begin{itemize}
    \item Les variables d'ecart $(x_{n+1}, x_{n+2}, \ldots, x_{n+m})$ sont de base
    \item Les variables de decision $(x_1, x_2, \ldots, x_n)$ sont hors-base
\end{itemize}

La solution associee au dictionnaire initial est :
\begin{align*}
x_1 = x_2 = \cdots = x_n &= 0\\
x_{n+i} = b_i &\quad \text{pour } i = 1, 2, \ldots, m
\end{align*}

L'ensemble des valeurs nulles pour les variables de decision est parfois appele l'\textcolor{defvert}{\textbf{origine}}.
\end{definition}

\begin{important}
\textbf{Proposition 3 :}

Le dictionnaire initial est realisable \textcolor{importantrouge}{\textbf{si et seulement si}} chaque membre droit $b_i$ est non negatif.
\end{important}

\begin{definition}
\textbf{Origine realisable :}

Les problemes de PL avec chaque membre droit $b_i$ non negatif sont appeles des problemes avec \textcolor{titrebleu}{\textbf{origine realisable}}.
\end{definition}

\begin{notecours}
\textbf{Remarque pratique :}

Dans ce cours, nous supposons generalement que nous avons une origine realisable (c'est-a-dire $b_i \geq 0$ pour tout $i$).

Si ce n'est pas le cas, des techniques speciales (comme la \textit{methode du simplexe a deux phases}) sont necessaires.
\end{notecours}

\newpage
\section{Strategie de la Methode du Simplexe}

\needspace{10\baselineskip}
\subsection{Idee Principale}

\begin{important}
\textbf{Strategie generale :}

La methode du simplexe travaille \textcolor{importantrouge}{\textbf{exclusivement}} avec des solutions de base realisables (c'est-a-dire des dictionnaires realisables) et ignore toutes les autres solutions realisables.

La grande strategie de la methode du simplexe est celle des \textcolor{defvert}{\textbf{ameliorations successives}} :

\begin{enumerate}
    \item Ayant trouve une solution de base realisable $x_1, x_2, \ldots, x_{n+m}$
    
    \item Essayer de passer a une autre solution de base realisable $\bar{x}_1, \bar{x}_2, \ldots, \bar{x}_{n+m}$ qui est \textcolor{titrebleu}{\textbf{meilleure}} au sens ou :
    $$\sum_{j=1}^{n} c_j\bar{x}_j > \sum_{j=1}^{n} c_jx_j$$
\end{enumerate}
\end{important}

\begin{astuce}
\textbf{En termes de dictionnaires :}

\begin{enumerate}
    \item On part d'un dictionnaire realisable
    \item On transforme ce dictionnaire en un nouveau dictionnaire realisable
    \item Le nouveau dictionnaire doit donner une valeur de $z$ \textcolor{defvert}{\textbf{strictement superieure}}
    \item On repete jusqu'a ne plus pouvoir ameliorer
\end{enumerate}
\end{astuce}

\newpage
\section{Un Deuxieme Exemple Detaille}

\needspace{10\baselineskip}
\subsection{Enonce du Probleme}

\begin{exemple}
\textbf{Programme lineaire sous forme standard :}

\begin{align*}
\text{maximiser} \quad & z = 5x_1 + 5x_2 + 3x_3\\
\text{sujet a} \quad & x_1 + 3x_2 + x_3 \leq 3\\
& -x_1 + 3x_3 \leq 2\\
& 2x_1 - x_2 + 2x_3 \leq 4\\
& 2x_1 + 3x_2 - x_3 \leq 2\\
& x_1, x_2, x_3 \geq 0
\end{align*}

\textbf{Dictionnaire initial realisable :}

\begin{align*}
x_4 &= 3 - x_1 - 3x_2 - x_3\\
x_5 &= 2 + x_1 - 3x_3\\
x_6 &= 4 - 2x_1 + x_2 - 2x_3\\
x_7 &= 2 - 2x_1 - 3x_2 + x_3\\
z &= 0 + 5x_1 + 5x_2 + 3x_3
\end{align*}

Base initiale : $B = \{4, 5, 6, 7\}$ (variables d'ecart)
\end{exemple}

\newpage
\subsection{Premiere Iteration : Amelioration}

\begin{methode}
\textbf{Comment ameliorer la solution ?}

Dictionnaire actuel :
\begin{align*}
x_4 &= 3 - x_1 - 3x_2 - x_3\\
x_5 &= 2 + x_1 - 3x_3\\
x_6 &= 4 - 2x_1 + x_2 - 2x_3\\
x_7 &= 2 - 2x_1 - 3x_2 + x_3\\
z &= 0 + 5x_1 + 5x_2 + 3x_3
\end{align*}

\textbf{Solution actuelle :} $x_1 = x_2 = x_3 = 0$, $z = 0$

\textbf{Idee :} Chercher une solution realisable qui donne une valeur de $z$ plus elevee.

\textbf{Strategie :}
\begin{itemize}
    \item Garder $x_2 = 0$ et $x_3 = 0$
    \item \textcolor{defvert}{\textbf{Augmenter}} la valeur de $x_1$
\end{itemize}

\textbf{Effet :} $z = 5x_1 > 0$ (amelioration !)
\end{methode}

\begin{exemple}
\textbf{Tests de valeurs :}

Si $x_1 = 1$ (en gardant $x_2 = x_3 = 0$) :
\begin{itemize}
    \item $x_4 = 3 - 1 = 2 \geq 0$ \checkmark
    \item $x_5 = 2 + 1 = 3 \geq 0$ \checkmark
    \item $x_6 = 4 - 2 = 2 \geq 0$ \checkmark
    \item $x_7 = 2 - 2 = 0 \geq 0$ \checkmark
    \item $z = 5$ \textcolor{defvert}{(realisable et meilleur !)}
\end{itemize}

Si $x_1 = 2$ (en gardant $x_2 = x_3 = 0$) :
\begin{itemize}
    \item $x_4 = 3 - 2 = 1 \geq 0$ \checkmark
    \item $x_5 = 2 + 2 = 4 \geq 0$ \checkmark
    \item $x_6 = 4 - 4 = 0 \geq 0$ \checkmark
    \item $x_7 = 2 - 4 = -2 < 0$ \textcolor{importantrouge}{\ding{55} (non realisable !)}
    \item $z = 10$
\end{itemize}

\textbf{Conclusion :} On ne peut pas augmenter $x_1$ trop.
\end{exemple}

\newpage
\subsection{Determination de la Borne Superieure}

\begin{methode}
\textbf{Question :} De combien peut-on augmenter $x_1$ (en gardant $x_2 = x_3 = 0$) tout en maintenant la realisabilite ?

\textbf{Analyse des contraintes :}

\begin{align*}
x_4 = 3 - x_1 \geq 0 &\Rightarrow x_1 \leq 3\\
x_5 = 2 + x_1 \geq 0 &\Rightarrow x_1 \geq -2 \text{ (toujours satisfait)}\\
x_6 = 4 - 2x_1 \geq 0 &\Rightarrow x_1 \leq 2\\
x_7 = 2 - 2x_1 \geq 0 &\Rightarrow x_1 \leq 1
\end{align*}

\textbf{Contrainte la plus stringente :} $x_1 \leq 1$

\textbf{Donc :} On peut augmenter $x_1$ jusqu'a $x_1 = 1$.
\end{methode}

\begin{important}
\textbf{Nouvelle solution :}

En posant $x_1 = 1$ (et $x_2 = x_3 = 0$) :
\begin{align*}
x_1 &= 1\\
x_2 &= 0\\
x_3 &= 0\\
x_4 &= 2\\
x_5 &= 3\\
x_6 &= 2\\
x_7 &= 0\\
z &= 5
\end{align*}

\textbf{Observations :}
\begin{itemize}
    \item $x_1 = 1 > 0$ : $x_1$ est devenue une \textcolor{titrebleu}{\textbf{variable de base}}
    \item $x_7 = 0$ : $x_7$ est devenue une \textcolor{defvert}{\textbf{variable hors-base}}
\end{itemize}
\end{important}

\newpage
\subsection{Construction du Nouveau Dictionnaire}

\begin{methode}
\textbf{Pivotage :}

On veut deplacer $x_1$ vers le membre gauche et $x_7$ vers le membre droit.

\textbf{Etape 1 :} Exprimer $x_1$ en fonction de $x_7$ (et des autres variables hors-base)

De la quatrieme equation du dictionnaire initial :
$$x_7 = 2 - 2x_1 - 3x_2 + x_3$$

On resout pour $x_1$ :
$$2x_1 = 2 - 3x_2 + x_3 - x_7$$
$$x_1 = 1 - \frac{3}{2}x_2 + \frac{1}{2}x_3 - \frac{1}{2}x_7$$

\textbf{Etape 2 :} Substituer cette expression dans les trois premieres equations et dans $z$

\textbf{Nouveau dictionnaire :}
\begin{align*}
x_1 &= 1 - \frac{3}{2}x_2 + \frac{1}{2}x_3 - \frac{1}{2}x_7\\
x_4 &= 2 - \frac{3}{2}x_2 - \frac{3}{2}x_3 + \frac{1}{2}x_7\\
x_5 &= 3 - \frac{3}{2}x_2 - \frac{5}{2}x_3 - \frac{1}{2}x_7\\
x_6 &= 2 + 4x_2 - 3x_3 + x_7\\
z &= 5 - \frac{5}{2}x_2 + \frac{11}{2}x_3 - \frac{5}{2}x_7
\end{align*}

\textbf{Nouvelle base :} $B = \{1, 4, 5, 6\}$
\end{methode}

\begin{astuce}
\textbf{Verification :}

En posant $x_2 = x_3 = x_7 = 0$ :
\begin{itemize}
    \item $x_1 = 1, x_4 = 2, x_5 = 3, x_6 = 2$ (toutes $\geq 0$) \checkmark
    \item $z = 5$ \checkmark
\end{itemize}

Le nouveau dictionnaire est \textcolor{defvert}{\textbf{realisable}} et donne une \textcolor{titrebleu}{\textbf{meilleure}} valeur de $z$ !
\end{astuce}

\newpage
\section{Une Iteration de la Methode du Simplexe}

\needspace{10\baselineskip}
\subsection{Structure Generale}

\begin{important}
\textbf{Iteration de la methode du simplexe :}

A chaque iteration, nous essayons d'augmenter la valeur de $z$ en rendant :
\begin{itemize}
    \item Une des variables du membre droit \textcolor{defvert}{\textbf{positive}} (variable entrante)
    \item Une des variables du membre gauche \textcolor{importantrouge}{\textbf{nulle}} (variable sortante)
\end{itemize}

tout en gardant les autres variables non negatives.

\textbf{En d'autres termes :}

Etant donne un dictionnaire realisable :
\begin{enumerate}
    \item Selectionner une \textcolor{defvert}{\textbf{variable entrante}}
    \item Trouver une \textcolor{importantrouge}{\textbf{variable sortante}}
    \item Construire le dictionnaire suivant par \textcolor{algoviolet}{\textbf{pivotage}}
\end{enumerate}
\end{important}

\newpage
\subsection{Choix de la Variable Entrante}

\begin{definition}
\textbf{Variable entrante :}

Une variable hors-base qui, lorsqu'elle devient positive, augmente la valeur de $z$, choisie pour entrer dans la base, est appelee la \textcolor{defvert}{\textbf{variable entrante}}.
\end{definition}

\begin{methode}
\textbf{Selection de la variable entrante :}

La variable entrante est une variable hors-base $x_j$ avec un coefficient \textcolor{titrebleu}{\textbf{positif}} $\bar{c}_j$ dans la derniere ligne du dictionnaire actuel :

$$z = z^* + \sum_{j \in N} \bar{c}_j x_j$$

ou $N$ est l'ensemble des indices des variables hors-base.

\textbf{Regles de choix :}
\begin{itemize}
    \item S'il y a plus d'un candidat pour entrer dans la base, \textcolor{defvert}{\textbf{n'importe lequel}} peut servir
    \item \textcolor{algoviolet}{\textbf{Pratique courante :}} choisir la variable ayant le plus grand coefficient $\bar{c}_j$
    \item Si $\bar{c}_j \leq 0$ pour tout $j \in N$, alors la solution actuelle est \textcolor{importantrouge}{\textbf{optimale}}
\end{itemize}
\end{methode}

\begin{astuce}
\textbf{Pourquoi le dictionnaire est optimal si tous les $\bar{c}_j \leq 0$ ?}

Car :
\begin{itemize}
    \item Toute solution realisable $x_1, x_2, \ldots, x_{n+m}$ satisfait $x_j \geq 0$ pour tout $j$
    \item Les variables hors-base sont fixees a zero dans la solution de base actuelle
    \item Si $\bar{c}_j \leq 0$ pour tout $j \in N$, augmenter une variable hors-base \textcolor{importantrouge}{\textbf{diminuerait}} $z$
    \item Donc la solution actuelle est optimale
\end{itemize}
\end{astuce}

\newpage
\subsection{Choix de la Variable Sortante}

\begin{definition}
\textbf{Variable sortante :}

La variable de base dont la non-negativite impose la borne superieure la plus stringente sur l'augmentation de la variable entrante est appelee la \textcolor{importantrouge}{\textbf{variable sortante}}.
\end{definition}

\begin{methode}
\textbf{Selection de la variable sortante :}

Supposons que $x_e$ soit la variable entrante (avec coefficient $\bar{c}_e > 0$).

Pour chaque variable de base $x_i$ (avec $i \in B$), si le coefficient de $x_e$ dans l'equation de $x_i$ est negatif (disons $-\bar{a}_{ie}$ avec $\bar{a}_{ie} > 0$), alors :

$$x_i = \bar{b}_i - \bar{a}_{ie}x_e - \cdots \geq 0$$

implique :
$$x_e \leq \frac{\bar{b}_i}{\bar{a}_{ie}}$$

\textbf{La variable sortante est celle qui donne le \textcolor{importantrouge}{plus petit} ratio} :

$$\min_{i \in B, \bar{a}_{ie} > 0} \frac{\bar{b}_i}{\bar{a}_{ie}}$$

\textbf{Regles supplementaires :}
\begin{itemize}
    \item S'il y a plus d'un candidat pour sortir, \textcolor{defvert}{\textbf{n'importe lequel}} peut servir
    \item S'il n'y a \textcolor{importantrouge}{\textbf{aucun candidat}} (tous les coefficients $\bar{a}_{ie} \leq 0$), le probleme est \textcolor{importantrouge}{\textbf{non borne}}
\end{itemize}
\end{methode}

\begin{important}
\textbf{Probleme non borne :}

S'il n'y a aucun candidat pour sortir de la base :
\begin{itemize}
    \item On peut rendre la variable entrante \textcolor{titrebleu}{\textbf{aussi grande qu'on veut}}
    \item La solution obtenue reste \textcolor{defvert}{\textbf{realisable}}
    \item Aucune valeur de variable de base n'a diminue
    \item La valeur de $z$ a augmente
    \item Plus la variable entrante augmente, plus $z$ augmente
    \item Comme la variable entrante n'a pas de borne superieure, $z$ n'a pas de borne superieure
    \item Le probleme de PL est \textcolor{importantrouge}{\textbf{non borne}}
\end{itemize}
\end{important}

\newpage
\subsection{Pivotage}

\begin{definition}
\textbf{Pivotage :}

Le nouveau dictionnaire est obtenu en :
\begin{enumerate}
    \item Faisant passer la variable entrante du membre droit au membre gauche, tandis que la variable sortante passe dans la direction opposee (sur la \textcolor{algoviolet}{\textbf{ligne du pivot}})
    
    \item Substituant la variable entrante de la ligne du pivot dans les lignes restantes du dictionnaire
\end{enumerate}

Le processus de calcul pour construire le nouveau dictionnaire est appele \textcolor{algoviolet}{\textbf{pivotage}}.
\end{definition}

\begin{methode}
\textbf{Etapes du pivotage :}

Supposons que $x_e$ entre et $x_s$ sorte.

\textbf{Etape 1 : Ligne du pivot}

Trouver l'equation de $x_s$ dans le dictionnaire actuel :
$$x_s = \bar{b}_s - \bar{a}_{se}x_e - \sum_{j \in N \setminus \{e\}} \bar{a}_{sj}x_j$$

Resoudre pour $x_e$ :
$$x_e = \frac{\bar{b}_s}{\bar{a}_{se}} - \frac{1}{\bar{a}_{se}}x_s - \sum_{j \in N \setminus \{e\}} \frac{\bar{a}_{sj}}{\bar{a}_{se}}x_j$$

\textbf{Etape 2 : Autres lignes}

Pour chaque autre variable de base $x_i$ ($i \in B \setminus \{s\}$), substituer l'expression de $x_e$ :
$$x_i = \bar{b}_i - \bar{a}_{ie}x_e - \cdots$$

devient :
$$x_i = \left(\bar{b}_i - \bar{a}_{ie} \cdot \frac{\bar{b}_s}{\bar{a}_{se}}\right) - \left(-\bar{a}_{ie} \cdot \frac{1}{\bar{a}_{se}}\right)x_s - \cdots$$

\textbf{Etape 3 : Fonction objectif}

Substituer $x_e$ dans l'expression de $z$ de maniere similaire.
\end{methode}

\begin{important}
\textbf{Proprietes apres pivotage :}

Une fois le pivotage termine :
\begin{itemize}
    \item Les $m$ variables de base doivent avoir des valeurs \textcolor{defvert}{\textbf{non negatives}}
    \item La valeur de la fonction objectif ne peut pas avoir \textcolor{titrebleu}{\textbf{diminue}}
\end{itemize}
\end{important}

\newpage
\section{Exemple Complet : Suite des Iterations}

\needspace{10\baselineskip}
\subsection{Deuxieme Iteration}

\begin{exemple}
\textbf{Dictionnaire ameliore (apres premiere iteration) :}

\begin{align*}
x_1 &= 1 - \frac{3}{2}x_2 + \frac{1}{2}x_3 - \frac{1}{2}x_7\\
x_4 &= 2 - \frac{3}{2}x_2 - \frac{3}{2}x_3 + \frac{1}{2}x_7\\
x_5 &= 3 - \frac{3}{2}x_2 - \frac{5}{2}x_3 - \frac{1}{2}x_7\\
x_6 &= 2 + 4x_2 - 3x_3 + x_7\\
z &= 5 - \frac{5}{2}x_2 + \frac{11}{2}x_3 - \frac{5}{2}x_7
\end{align*}

\textbf{Solution actuelle :} $x_1 = 1, x_2 = 0, x_3 = 0, x_4 = 2, x_5 = 3, x_6 = 2, x_7 = 0$, $z = 5$
\end{exemple}

\begin{methode}
\textbf{Selection de la variable entrante :}

Dans l'equation de $z$ : $z = 5 - \frac{5}{2}x_2 + \frac{11}{2}x_3 - \frac{5}{2}x_7$

\textbf{Coefficients des variables hors-base :}
\begin{itemize}
    \item $\bar{c}_2 = -\frac{5}{2} < 0$ (ne peut pas entrer)
    \item $\bar{c}_3 = \frac{11}{2} > 0$ \textcolor{defvert}{(peut entrer)}
    \item $\bar{c}_7 = -\frac{5}{2} < 0$ (ne peut pas entrer)
\end{itemize}

\textbf{Variable entrante :} $x_3$ (seule variable avec coefficient positif)
\end{methode}

\begin{methode}
\textbf{Selection de la variable sortante :}

Coefficient de $x_3$ dans chaque equation :

\begin{align*}
x_1 &: +\frac{1}{2} \quad (\text{positif, pas de borne})\\
x_4 &: -\frac{3}{2} \quad \Rightarrow x_3 \leq \frac{2}{3/2} = \frac{4}{3}\\
x_5 &: -\frac{5}{2} \quad \Rightarrow x_3 \leq \frac{3}{5/2} = \frac{6}{5}\\
x_6 &: -3 \quad \Rightarrow x_3 \leq \frac{2}{3}
\end{align*}

\textbf{Borne la plus stringente :} $x_3 \leq \frac{2}{3}$ (imposee par $x_6$)

\textbf{Variable sortante :} $x_6$
\end{methode}

\newpage
\subsection{Troisieme Dictionnaire}

\begin{methode}
\textbf{Pivotage :}

De l'equation de $x_6$ :
$$x_6 = 2 + 4x_2 - 3x_3 + x_7$$

On resout pour $x_3$ :
$$3x_3 = 2 + 4x_2 + x_7 - x_6$$
$$x_3 = \frac{2}{3} + \frac{4}{3}x_2 + \frac{1}{3}x_7 - \frac{1}{3}x_6$$

Substitution dans les autres equations...
\end{methode}

\begin{exemple}
\textbf{Troisieme dictionnaire :}

\begin{align*}
x_3 &= \frac{2}{3} + \frac{4}{3}x_2 + \frac{1}{3}x_7 - \frac{1}{3}x_6\\
x_1 &= \frac{4}{3} - \frac{5}{6}x_2 - \frac{1}{3}x_7 - \frac{1}{6}x_6\\
x_4 &= 1 - \frac{7}{2}x_2 + \frac{1}{2}x_6\\
x_5 &= \frac{4}{3} - \frac{29}{6}x_2 - \frac{4}{3}x_7 + \frac{5}{6}x_6\\
z &= \frac{26}{3} + \frac{29}{6}x_2 - \frac{2}{3}x_7 - \frac{11}{6}x_6
\end{align*}

\textbf{Solution amelioree :} 
$$x_1 = \frac{4}{3}, x_2 = 0, x_3 = \frac{2}{3}, x_4 = 1, x_5 = \frac{4}{3}, x_6 = 0, x_7 = 0$$
$$z = \frac{26}{3} \approx 8.67$$
\end{exemple}

\newpage
\subsection{Troisieme Iteration}

\begin{methode}
\textbf{Variable entrante :}

Dans $z = \frac{26}{3} + \frac{29}{6}x_2 - \frac{2}{3}x_7 - \frac{11}{6}x_6$

Seul coefficient positif : $\bar{c}_2 = \frac{29}{6} > 0$

\textbf{Variable entrante :} $x_2$
\end{methode}

\begin{methode}
\textbf{Variable sortante :}

Ratios :
\begin{align*}
x_3 &: +\frac{4}{3} \quad (\text{pas de borne})\\
x_1 &: -\frac{5}{6} \quad \Rightarrow x_2 \leq \frac{4/3}{5/6} = \frac{8}{5}\\
x_4 &: -\frac{7}{2} \quad \Rightarrow x_2 \leq \frac{1}{7/2} = \frac{2}{7}\\
x_5 &: -\frac{29}{6} \quad \Rightarrow x_2 \leq \frac{4/3}{29/6} = \frac{8}{29}
\end{align*}

\textbf{Minimum :} $\frac{8}{29}$ (impose par $x_5$)

\textbf{Variable sortante :} $x_5$
\end{methode}

\newpage
\subsection{Dictionnaire Final (Optimal)}

\begin{exemple}
\textbf{Quatrieme dictionnaire (apres pivotage) :}

\begin{align*}
x_2 &= \frac{8}{29} - \frac{8}{29}x_7 + \frac{5}{29}x_6 - \frac{6}{29}x_5\\
x_3 &= \frac{30}{29} - \frac{1}{29}x_7 - \frac{3}{29}x_6 - \frac{8}{29}x_5\\
x_1 &= \frac{32}{29} - \frac{3}{29}x_7 - \frac{9}{29}x_6 + \frac{5}{29}x_5\\
x_4 &= \frac{1}{29} + \frac{28}{29}x_7 - \frac{3}{29}x_6 + \frac{21}{29}x_5\\
z &= 10 - 2x_7 - x_6 - x_5
\end{align*}

\textbf{Solution optimale :}
$$x_1 = \frac{32}{29}, \quad x_2 = \frac{8}{29}, \quad x_3 = \frac{30}{29}$$
$$x_4 = \frac{1}{29}, \quad x_5 = 0, \quad x_6 = 0, \quad x_7 = 0$$

\textbf{Valeur optimale :} $z = 10$
\end{exemple}

\begin{important}
\textbf{Critere d'optimalite :}

Dans l'equation de $z$ : $z = 10 - 2x_7 - x_6 - x_5$

\textcolor{importantrouge}{\textbf{Tous les coefficients sont negatifs ou nuls !}}

Aucune variable hors-base ne peut entrer dans la base sans faire diminuer $z$.

\textcolor{defvert}{\textbf{Le dictionnaire actuel decrit une solution optimale.}}
\end{important}

\newpage
\section{Critere d'Arret et Solutions Optimales}

\needspace{10\baselineskip}
\subsection{Critere d'Arret}

\begin{definition}
\textbf{Critere d'arret de la methode du simplexe :}

Repeter le processus jusqu'a ce qu'aucune variable hors-base ne puisse entrer dans la base sans faire diminuer la valeur de $z$.

La solution realisable associee au dictionnaire actuel est une \textcolor{defvert}{\textbf{solution optimale}}.

\textbf{Formellement :} Si $\bar{c}_j \leq 0$ pour tout $j \in N$ (variables hors-base), alors \textcolor{importantrouge}{\textbf{STOP}} - solution optimale trouvee.
\end{definition}

\begin{astuce}
\textbf{Resume de l'algorithme :}

\begin{enumerate}
    \item \textbf{Initialisation :} Partir d'un dictionnaire realisable (souvent le dictionnaire initial si origine realisable)
    
    \item \textbf{Test d'optimalite :} Si tous les coefficients $\bar{c}_j \leq 0$ dans l'equation de $z$, \textcolor{defvert}{STOP} (optimal)
    
    \item \textbf{Choix variable entrante :} Choisir $x_e$ avec $\bar{c}_e > 0$
    
    \item \textbf{Test de non-borne :} Si aucun coefficient $\bar{a}_{ie} > 0$ pour $i \in B$, \textcolor{importantrouge}{STOP} (non borne)
    
    \item \textbf{Choix variable sortante :} Calculer $\min_{i \in B, \bar{a}_{ie} > 0} \frac{\bar{b}_i}{\bar{a}_{ie}}$ pour determiner $x_s$
    
    \item \textbf{Pivotage :} Construire le nouveau dictionnaire
    
    \item \textbf{Retour a l'etape 2}
\end{enumerate}
\end{astuce}

\newpage
\subsection{Unicite de la Solution Optimale}

\begin{exemple}
\textbf{Exemple avec solution optimale unique :}

Rappelons notre premier exemple :

\textbf{Dictionnaire final :}
\begin{align*}
x_3 &= 1 + x_2 + 3x_4 - 2x_6\\
x_1 &= 2 - 2x_2 - 2x_4 + x_6\\
x_5 &= 1 + 5x_2 + 2x_4\\
z &= 13 - 3x_2 - x_4 - x_6
\end{align*}

La derniere ligne montre que toute solution realisable avec $z = 13$ doit satisfaire :
$$x_2 = x_4 = x_6 = 0$$

Le reste du dictionnaire montre que toute telle solution doit satisfaire :
$$x_1 = 2, \quad x_3 = 1, \quad x_5 = 1$$

\textcolor{defvert}{\textbf{Il n'y a donc qu'une seule solution optimale.}}
\end{exemple}

\newpage
\subsection{Solutions Optimales Multiples}

\begin{definition}
\textbf{Detection de solutions optimales multiples :}

Considerons le dictionnaire suivant :

\begin{align*}
x_4 &= 3 + x_2 - 2x_5 + 7x_3\\
x_1 &= 1 - 5x_2 + 6x_5 - 8x_3\\
x_6 &= 4 + 9x_2 + 2x_5 - x_3\\
z &= 8 - x_3
\end{align*}

La derniere ligne montre que toute solution optimale satisfait $x_3 = 0$, mais \textcolor{importantrouge}{\textbf{pas necessairement}} $x_2 = 0$ et $x_5 = 0$.
\end{definition}

\begin{exemple}
\textbf{Caracterisation des solutions optimales :}

Pour les solutions avec $x_3 = 0$, le reste du dictionnaire implique :

\begin{align*}
x_4 &= 3 + x_2 - 2x_5\\
x_1 &= 1 - 5x_2 + 6x_5\\
x_6 &= 4 + 9x_2 + 2x_5
\end{align*}

Toute solution optimale provient de valeurs de $x_2$ et $x_5$ telles que :

\begin{align*}
-x_2 + 2x_5 &\leq 3\\
5x_2 - 6x_5 &\leq 1\\
-9x_2 - 2x_5 &\leq 4\\
x_2, x_5 &\geq 0
\end{align*}

\textcolor{titrebleu}{\textbf{Il existe une infinite de solutions optimales !}}

Ce sont tous les points du polyedre defini ci-dessus.
\end{exemple}

\begin{astuce}
\textbf{Comment detecter des solutions optimales multiples ?}

Dans un dictionnaire optimal, si une variable hors-base a un coefficient \textcolor{defvert}{\textbf{nul}} (pas strictement negatif) dans l'equation de $z$, alors il existe \textcolor{titrebleu}{\textbf{plusieurs solutions optimales}}.

Dans l'exemple ci-dessus : $\bar{c}_2 = 0$ et $\bar{c}_5 = 0$.
\end{astuce}

\newpage
\section{Format Tableau}

\needspace{10\baselineskip}
\subsection{Introduction au Tableau}

\begin{notecours}
\textbf{Une autre facon de presenter la methode du simplexe :}

Le format \textcolor{algoviolet}{\textbf{tableau}} est probablement plus populaire, bien que moins direct que le format dictionnaire.

Il consiste a enregistrer uniquement les \textcolor{titrebleu}{\textbf{coefficients}} des variables, ainsi que les membres droits, sous forme de tableau (matrice augmentee).
\end{notecours}

\begin{exemple}
\textbf{Dictionnaire (forme modifiee) pour notre premier exemple :}

\begin{align*}
2x_1 + 3x_2 + x_3 + x_4 &= 5\\
4x_1 + x_2 + 2x_3 + x_5 &= 11\\
3x_1 + 4x_2 + 2x_3 + x_6 &= 8\\
-z + 5x_1 + 4x_2 + 3x_3 &= 0
\end{align*}

\textbf{Tableau correspondant :}

\begin{center}
\begin{tabular}{ccccccc|c}
\toprule
$x_1$ & $x_2$ & $x_3$ & $x_4$ & $x_5$ & $x_6$ & $z$ & RHS \\
\midrule
2 & 3 & 1 & 1 & 0 & 0 & 0 & 5 \\
4 & 1 & 2 & 0 & 1 & 0 & 0 & 11 \\
3 & 4 & 2 & 0 & 0 & 1 & 0 & 8 \\
\midrule
5 & 4 & 3 & 0 & 0 & 0 & -1 & 0 \\
\bottomrule
\end{tabular}
\end{center}

RHS = Right-Hand Side (membre droit)
\end{exemple}

\newpage
\subsection{Operations sur le Tableau}

\begin{methode}
\textbf{Selection de la variable entrante :}

\begin{center}
\begin{tabular}{ccccccc|c}
\toprule
$x_1$ & $x_2$ & $x_3$ & $x_4$ & $x_5$ & $x_6$ & $z$ & RHS \\
\midrule
2 & 3 & 1 & 1 & 0 & 0 & 0 & 5 \\
4 & 1 & 2 & 0 & 1 & 0 & 0 & 11 \\
3 & 4 & 2 & 0 & 0 & 1 & 0 & 8 \\
\midrule
\textcolor{defvert}{5} & \textcolor{defvert}{4} & \textcolor{defvert}{3} & 0 & 0 & 0 & -1 & 0 \\
\bottomrule
\end{tabular}
\end{center}

\textbf{Regles :}
\begin{itemize}
    \item Examiner tous les nombres (sauf le plus a droite) dans la \textcolor{algoviolet}{\textbf{derniere ligne}}
    \item S'ils sont \textcolor{defvert}{\textbf{tous non positifs}} : le tableau decrit une solution optimale
    \item Sinon : choisir une colonne avec un nombre \textcolor{titrebleu}{\textbf{positif}} dans la derniere ligne (\textcolor{importantrouge}{\textbf{colonne pivot}} $\to$ variable entrante)
\end{itemize}

\textbf{Choix :} Colonne de $x_1$ (coefficient = 5)
\end{methode}

\begin{methode}
\textbf{Selection de la variable sortante :}

\begin{center}
\begin{tabular}{ccccccc|c|c}
\toprule
$x_1$ & $x_2$ & $x_3$ & $x_4$ & $x_5$ & $x_6$ & $z$ & RHS & Ratio \\
\midrule
\textcolor{defvert}{2} & 3 & 1 & 1 & 0 & 0 & 0 & 5 & $5/2 = 2.5$ \\
\textcolor{defvert}{4} & 1 & 2 & 0 & 1 & 0 & 0 & 11 & $11/4 = 2.75$ \\
\textcolor{defvert}{3} & 4 & 2 & 0 & 0 & 1 & 0 & 8 & $8/3 \approx 2.67$ \\
\midrule
5 & 4 & 3 & 0 & 0 & 0 & -1 & 0 & - \\
\bottomrule
\end{tabular}
\end{center}

\textbf{Regles :}
\begin{itemize}
    \item Pour chaque ligne dont l'entree $r$ dans la colonne pivot est \textcolor{titrebleu}{\textbf{positive}}, calculer le ratio $s/r$ ou $s$ est l'entree dans la colonne la plus a droite
    \item Si toutes les entrees de la colonne pivot sont non positives : le probleme de PL est \textcolor{importantrouge}{\textbf{non borne}}
    \item La ligne avec le \textcolor{defvert}{\textbf{plus petit ratio}} est la \textcolor{importantrouge}{\textbf{ligne pivot}} ($\to$ variable sortante)
\end{itemize}

\textbf{Choix :} Ligne 1 (ratio = 2.5, le plus petit)

\textbf{Element pivot :} L'entree a l'intersection de la colonne pivot et de la ligne pivot (\circled{2})
\end{methode}

\newpage
\subsection{Pivotage sur le Tableau}

\begin{methode}
\textbf{Operations de pivotage :}

\textbf{Tableau initial :}
\begin{center}
\begin{tabular}{ccccccc|c}
\toprule
$x_1$ & $x_2$ & $x_3$ & $x_4$ & $x_5$ & $x_6$ & $z$ & RHS \\
\midrule
\circled{2} & 3 & 1 & 1 & 0 & 0 & 0 & 5 \\
4 & 1 & 2 & 0 & 1 & 0 & 0 & 11 \\
3 & 4 & 2 & 0 & 0 & 1 & 0 & 8 \\
\midrule
5 & 4 & 3 & 0 & 0 & 0 & -1 & 0 \\
\bottomrule
\end{tabular}
\end{center}

\textbf{Etape 1 :} Diviser chaque entree de la ligne pivot par l'element pivot

$$\text{Ligne 1 nouvelle} = \text{Ligne 1 ancienne} / 2$$

\begin{center}
\begin{tabular}{ccccccc|c}
\toprule
$x_1$ & $x_2$ & $x_3$ & $x_4$ & $x_5$ & $x_6$ & $z$ & RHS \\
\midrule
\textcolor{defvert}{1} & \textcolor{defvert}{3/2} & \textcolor{defvert}{1/2} & \textcolor{defvert}{1/2} & 0 & 0 & 0 & \textcolor{defvert}{5/2} \\
4 & 1 & 2 & 0 & 1 & 0 & 0 & 11 \\
3 & 4 & 2 & 0 & 0 & 1 & 0 & 8 \\
\midrule
5 & 4 & 3 & 0 & 0 & 0 & -1 & 0 \\
\bottomrule
\end{tabular}
\end{center}

\textbf{Etape 2 :} Pour chaque ligne restante, soustraire un multiple approprie de la nouvelle ligne pivot

\textbf{Objectif :} Rendre toutes les entrees de la colonne pivot (sauf l'element pivot) egales a zero

$$\text{Ligne 2 nouvelle} = \text{Ligne 2 ancienne} - 4 \times \text{Ligne 1 nouvelle}$$
$$\text{Ligne 3 nouvelle} = \text{Ligne 3 ancienne} - 3 \times \text{Ligne 1 nouvelle}$$
$$\text{Ligne 4 nouvelle} = \text{Ligne 4 ancienne} - 5 \times \text{Ligne 1 nouvelle}$$
\end{methode}

\begin{exemple}
\textbf{Tableau apres pivotage :}

\begin{center}
\begin{tabular}{ccccccc|c}
\toprule
$x_1$ & $x_2$ & $x_3$ & $x_4$ & $x_5$ & $x_6$ & $z$ & RHS \\
\midrule
1 & 3/2 & 1/2 & 1/2 & 0 & 0 & 0 & 5/2 \\
0 & -5 & 0 & -2 & 1 & 0 & 0 & 1 \\
0 & -1/2 & 1/2 & -3/2 & 0 & 1 & 0 & 1/2 \\
\midrule
0 & -7/2 & 1/2 & -5/2 & 0 & 0 & -1 & -25/2 \\
\bottomrule
\end{tabular}
\end{center}

\textbf{Interpretation :}
\begin{itemize}
    \item Variables de base : $x_1, x_5, x_6$ (colonnes avec un seul 1 et des 0 ailleurs)
    \item Variables hors-base : $x_2, x_3, x_4$
    \item Solution : $x_1 = 5/2$, $x_5 = 1$, $x_6 = 1/2$ ; autres variables = 0
    \item Valeur de $z$ : $25/2 = 12.5$ (attention au signe !)
\end{itemize}

\textbf{Note :} Dans la derniere ligne, $z$ a un coefficient de -1, donc la valeur de $z$ est $-(-25/2) = 25/2$.
\end{exemple}

\newpage
\section{Resume et Conseils Pratiques}

\needspace{10\baselineskip}
\subsection{Algorithme du Simplexe - Resume}

\begin{important}
\textbf{Algorithme du simplexe (version dictionnaire) :}

\textbf{Entree :} Probleme de PL sous forme standard avec origine realisable

\textbf{Sortie :} Solution optimale ou indication que le probleme est non borne

\textbf{Etapes :}
\begin{enumerate}
    \item \textcolor{titrebleu}{\textbf{Initialisation :}}
    \begin{itemize}
        \item Former le dictionnaire initial avec les variables d'ecart
        \item Verifier qu'il est realisable ($b_i \geq 0$ pour tout $i$)
    \end{itemize}
    
    \item \textcolor{defvert}{\textbf{Test d'optimalite :}}
    \begin{itemize}
        \item Si tous les coefficients dans l'equation de $z$ sont $\leq 0$ : \textbf{STOP} (optimal)
    \end{itemize}
    
    \item \textcolor{algoviolet}{\textbf{Choix de la variable entrante :}}
    \begin{itemize}
        \item Choisir une variable hors-base avec coefficient positif dans $z$
    \end{itemize}
    
    \item \textcolor{exempleorange}{\textbf{Test de borne :}}
    \begin{itemize}
        \item Si aucun coefficient de la variable entrante n'est positif : \textbf{STOP} (non borne)
    \end{itemize}
    
    \item \textcolor{importantrouge}{\textbf{Choix de la variable sortante :}}
    \begin{itemize}
        \item Calculer les ratios $\bar{b}_i / \bar{a}_{ie}$ pour $\bar{a}_{ie} > 0$
        \item Choisir le minimum
    \end{itemize}
    
    \item \textcolor{titrebleu}{\textbf{Pivotage :}}
    \begin{itemize}
        \item Construire le nouveau dictionnaire
        \item Retourner a l'etape 2
    \end{itemize}
\end{enumerate}
\end{important}

\newpage
\subsection{Conseils et Pieges Courants}

\begin{astuce}
\textbf{Conseils pratiques :}

\begin{enumerate}
    \item \textbf{Verification de la realisabilite :}
    \begin{itemize}
        \item Toujours verifier que les valeurs des variables de base sont $\geq 0$
        \item Si une valeur devient negative, il y a une erreur de calcul
    \end{itemize}
    
    \item \textbf{Amelioration de $z$ :}
    \begin{itemize}
        \item A chaque iteration, $z$ doit augmenter (ou au moins ne pas diminuer)
        \item Si $z$ diminue, il y a une erreur
    \end{itemize}
    
    \item \textbf{Choix de la variable entrante :}
    \begin{itemize}
        \item N'importe quelle variable avec coefficient positif peut entrer
        \item Le choix du plus grand coefficient (regle de Dantzig) est souvent efficace
        \item Mais d'autres regles existent (regle de Bland, etc.)
    \end{itemize}
    
    \item \textbf{Calcul des ratios :}
    \begin{itemize}
        \item Ne considerer que les coefficients \textcolor{importantrouge}{\textbf{strictement positifs}}
        \item Un coefficient nul ou negatif n'impose pas de borne
    \end{itemize}
    
    \item \textbf{Pivotage :}
    \begin{itemize}
        \item Bien identifier l'element pivot
        \item Verifier les calculs (facile de faire des erreurs arithmetiques)
        \item Garder les fractions si possible (plus precis que les decimales)
    \end{itemize}
\end{enumerate}
\end{astuce}

\begin{important}
\textbf{Pieges courants :}

\begin{itemize}
    \item \textcolor{importantrouge}{\textbf{Oublier le signe}} dans l'equation de $z$ (dans le format tableau, $z$ a un coefficient de -1)
    
    \item \textcolor{importantrouge}{\textbf{Confondre}} variables de base et hors-base
    
    \item \textcolor{importantrouge}{\textbf{Ne pas verifier}} la realisabilite du dictionnaire initial
    
    \item \textcolor{importantrouge}{\textbf{Calculer le mauvais ratio}} (numerateur/denominateur inverses)
    
    \item \textcolor{importantrouge}{\textbf{Ne pas tester}} le critere d'arret a chaque iteration
    
    \item \textcolor{importantrouge}{\textbf{Croire}} qu'il y a toujours une solution optimale unique
\end{itemize}
\end{important}

\newpage
\section{Exercices}

\begin{exemple}
\textbf{Exercice 1 : Application directe}

Resolvez par la methode du simplexe :

\begin{align*}
\text{maximiser} \quad & z = 3x_1 + 2x_2\\
\text{sujet a} \quad & x_1 + x_2 \leq 4\\
& 2x_1 + x_2 \leq 6\\
& x_1, x_2 \geq 0
\end{align*}

\textbf{Questions :}
\begin{enumerate}
    \item Ecrire le dictionnaire initial
    \item Effectuer les iterations du simplexe
    \item Donner la solution optimale et la valeur optimale
\end{enumerate}
\end{exemple}

\begin{exemple}
\textbf{Exercice 2 : Detection de cas particuliers}

Pour chacun des dictionnaires suivants, determiner :
\begin{itemize}
    \item S'il est realisable
    \item S'il est optimal
    \item S'il y a des solutions optimales multiples
    \item S'il indique un probleme non borne
\end{itemize}

\textbf{a)} 
\begin{align*}
x_3 &= 5 - 2x_1 + x_2\\
x_4 &= 3 + x_1 - x_2\\
z &= 10 + 2x_1 - 3x_2
\end{align*}

\textbf{b)} 
\begin{align*}
x_3 &= 2 + x_1 + x_2\\
x_4 &= 1 + 2x_1 + 3x_2\\
z &= 8 - x_1 - 2x_2
\end{align*}

\textbf{c)} 
\begin{align*}
x_1 &= 4 - x_3 + 2x_4\\
x_2 &= 3 + 2x_3 - x_4\\
z &= 15 - 0x_3 - x_4
\end{align*}
\end{exemple}

\newpage
\begin{exemple}
\textbf{Exercice 3 : Format tableau}

Convertissez le dictionnaire suivant en format tableau :

\begin{align*}
x_1 &= 3 - 2x_3 + x_5\\
x_2 &= 2 + x_3 - 2x_5\\
x_4 &= 1 - x_3 + x_5\\
z &= 7 - 3x_3 + 2x_5
\end{align*}

Puis effectuez une iteration du simplexe en utilisant le format tableau.
\end{exemple}

\begin{exemple}
\textbf{Exercice 4 : Probleme complet}

Une entreprise fabrique deux types de produits A et B. La production de A necessite 2 heures sur la machine M1 et 1 heure sur la machine M2. La production de B necessite 1 heure sur M1 et 2 heures sur M2. Les machines M1 et M2 sont disponibles respectivement 100 et 80 heures par semaine. Les profits unitaires sont de 30€ pour A et 40€ pour B.

\textbf{Questions :}
\begin{enumerate}
    \item Formuler ce probleme en PL
    \item Le resoudre par la methode du simplexe
    \item Interpreter la solution optimale
\end{enumerate}
\end{exemple}

\newpage
\section{Conclusion du Chapitre}

\begin{important}
\textbf{Points cles a retenir :}

\begin{enumerate}
    \item \textcolor{titrebleu}{\textbf{Variables d'ecart :}}
    \begin{itemize}
        \item Transforment les inegalites en egalites
        \item Permettent de travailler avec un systeme d'equations
    \end{itemize}
    
    \item \textcolor{defvert}{\textbf{Dictionnaires :}}
    \begin{itemize}
        \item Representent les solutions de base
        \item Variables de base (membre gauche) et hors-base (membre droit)
        \item Dictionnaire realisable : valeurs de base $\geq 0$
    \end{itemize}
    
    \item \textcolor{algoviolet}{\textbf{Methode du simplexe :}}
    \begin{itemize}
        \item Ameliorations successives de solutions de base
        \item Variable entrante (coefficient positif dans $z$)
        \item Variable sortante (ratio minimum)
        \item Pivotage pour obtenir nouveau dictionnaire
    \end{itemize}
    
    \item \textcolor{exempleorange}{\textbf{Criteres d'arret :}}
    \begin{itemize}
        \item Optimal : tous coefficients $\leq 0$ dans $z$
        \item Non borne : aucun coefficient positif pour variable entrante
    \end{itemize}
    
    \item \textcolor{importantrouge}{\textbf{Formats :}}
    \begin{itemize}
        \item Dictionnaire (plus intuitif)
        \item Tableau (plus compact, plus utilise en pratique)
    \end{itemize}
\end{enumerate}
\end{important}

\begin{astuce}
\textbf{Prochaines etapes :}

Dans les prochains chapitres, nous etudierons :
\begin{itemize}
    \item La methode du simplexe a deux phases (cas sans origine realisable)
    \item La theorie de la dualite
    \item Le simplexe revise
    \item L'analyse de sensibilite
    \item La degenerescence et le cyclage
\end{itemize}
\end{astuce}

\end{document}
